% SOURCE_AUTHORITY: sga5_fr_workpass.tex, lines 10611--10631 (LF-normalized unit).
% SOURCE_SHA256: 791F4EFFC5E02832D5D77ED03518C8156D6F07E4C8238B03545DB93D883FBB28
\emph{Demonstração.} Seja $j:X\hookrightarrow\mathbb P^r$ uma imersão fechada e seja $\xi_X$ (resp. $\xi_Y$) a classe em $H^2(X,\mathrm Q_\ell(1))$ (resp. $H^2(Y,\mathbb Q_\ell(1))$) de $\mathcal O_X(1)$ (resp. $\mathcal O_Y(1)$). Denota-se por $\alpha:Y\hookrightarrow X$ a inclusão canônica. Para $q\le n-1$, o diagrama
\[
\begin{array}{ccc}
H^q(X,L)&\xrightarrow{\alpha^*}&H^q(Y,L|Y)\\
\cup\,\xi_X^{\,n-1-q}\downarrow&&\downarrow\cup\,\xi_Y^{\,n-1-q}\\
H^{2n-2-q}(X,L(n-1-q))&\xrightarrow{\alpha^*}&
H^{2n-2-q}(Y,(L|Y)(n-1-q))
\end{array}
\]
é comutativo e a coluna da direita é (conjecturalmente) um isomorfismo. Ponhamos $p=2n-2-q$; para $q\le n-2$, ou seja, $p\ge n$, a linha superior é um isomorfismo (7.2), de onde se deduz a afirmação (i). A afirmação (ii) demonstra-se de igual modo considerando desta vez o diagrama seguinte, com $p\le n-1$:
\[
\begin{array}{ccc}
H^p(Y,(L|Y)(-1))&\xrightarrow{\mathrm{Gysin}}&H^{p+2}(X,L)\\
\cup\,\xi_Y^{\,n-1-p}\downarrow&&\downarrow\cup\,\xi_X^{\,n-1-p}\\
H^{2n-2-p}(Y,(L|Y)(n-2-p))&\xrightarrow{\mathrm{Gysin}}&
H^{2n-p}(X,L(n-1-p)).
\end{array}
\]
onde a linha inferior é um isomorfismo para $2n-2-p\ge n$, ou seja, $p\le n-2$ (7.2 (ii)).

Para uma discussão dessas conjecturas e de suas consequências, remete-se o leitor à exposição de Kleiman, \emph{Algebraic cycles and the Weil conjectures}, em «Dix exposés sur la cohomologie des schémas'').
